\section{Pipeline与Agent双轨：防止发现过程验证自己}

Pipeline 批量筛选单描述符并构造二/三元公式；Agent 每次只评估一个研究者明确指定的活动描述符。两条轨道共享冻结的 raw CSV、描述符注册表及修订号，但分别写入 \code{results/pipeline/} 与 \code{results/agent/}，在用户授权 C9 前不能读取对方结果。

这种隔离借鉴独立复核与三角验证思想。若 Pipeline 的组合搜索和 Agent 的单变量审计都指向 Na 网络连通性，说明该方向值得优先检查，但不能把一致性解释为因果证明。若两轨不一致，也不是简单失败：它可能揭示组合公式依赖特定协同、单变量信号不稳，或 Pipeline 的搜索空间与 Agent 的人工假设不一致。

\section{一个候选何时才算“可解释的局域结构信号”}

本框架将候选资格定义为一条可审计链，而不是单一阈值：

\begin{enumerate}
  \item 结构文件存在、可解析，描述符有明确单位、物理族和计算 provenance；
  \item 描述符与 $Y$ 的关系不是仅由 system/anion 均值差驱动；
  \item 在半样本 Lasso 中反复入选，并超过固定噪声基线；
  \item 若形成组合，运算来自预先声明的物理规则且量纲可解释；
  \item 在至少一种常规验证和 LOSO 中具有可用证据，方向冲突被显式报告；
  \item V1--V4 不显示明显随机公式、单体系驱动或极端重采样不确定性；
  \item 候选能够返回具体结构对象，例如 Na 配位、位点网络或间隙通道，并可提出可验证的材料学预测。
\end{enumerate}

只有完成这些步骤，候选才进入后续 BVSE/BVEL、NEB、AIMD 或实验验证。静态 CIF 描述符无法直接给出有限温度迁移势垒，也不能代表晶界、相变和制样工艺。因此，Pipeline 的最终产物是“优先验证列表”，不是材料机制终局。

\section{当前实现状态与本文能够支持的结论}

修复分支通过 70 项自动化测试、Python 全量编译和 CLI 冒烟检查。测试覆盖带电物种、混合占位、周期最小像、Wyckoff 多样性、全缺失失败关闭、折内插补、不可行 CV 的 skip 语义、Stage 1--2 合约、硬物理族容量、真实三元式生成、严格 JSON 往返、V1--V4 provenance、双轨隔离和冻结身份检查。

但当前 checkout 中 raw CSV 引用的 84 个 CIF 不可用。Pipeline 与 Agent 因而会在任何研究结果写入前失败。这一状态证明了失败关闭设计能够阻止全 NaN 特征被误写为材料学发现，却不能证明哪个真实描述符与电导率相关。因此，本稿不报告“最佳单描述符”“最佳组合”“组合提升百分比”或新的迁移机制。补齐并冻结 CIF 后，必须从 Stage 0 重新运行，并将真实数据图替换本文中的概念示意图。

\section{讨论：该框架的优势、代价与适用边界}

该框架的首要优势是把每个模型的职责限制在一个可解释问题内。Ridge 不是被笼统称为“机器学习模型”，而是先用于估计体系背景，后用于折外探测简单候选；Lasso 不用于解释系数大小，而只用于生成可重复的选择事件；Spearman 不承担复杂非线性拟合，而用于保留结构趋势方向；LOSO 不追求平均分最高，而是主动制造最困难的未知体系测试。这种职责分离减少了算法之间的语义混乱。

代价是框架较为保守。秩感知残差化可能移除一部分真实但与体系高度绑定的机制；稳定性选择可能在一组高度相关的同类描述符之间分散频率；受限公式搜索可能遗漏未写入规则表的真实表达；LOSO 在只有三个体系时方差很大。保守并不等于这些方法绝对优于 RF、XGBoost、GNN 或 SISSO，而是它们与当前目标和数据规模更匹配。未来样本扩展到数千并具有独立外部体系后，可以将图网络作为预测基线，或在已经完成去混杂与稳定压缩的候选空间上运行更广泛的符号回归。

最后，关联稳定性不等于机制稳定性。即使一个 Na 网络描述符通过所有统计门槛，也可能只是未记录缺陷浓度或相纯度的代理。真正的机制证据需要在结构扰动、组成系列、温度依赖电导、迁移势垒和有限温度轨迹中得到一致支持。本文框架的价值，是把海量结构变量压缩成少量、证据链透明且可以被计算和实验反驳的候选假设。
